\documentclass{amsart}
% For dealing with references we use the comment environment
\usepackage{verbatim}
\newenvironment{reference}{\comment}{\endcomment}
%\newenvironment{reference}{}{}
\newenvironment{slogan}{\comment}{\endcomment}
\newenvironment{history}{\comment}{\endcomment}

% For commutative diagrams we use Xy-pic
\usepackage[all]{xy}

% We use 2cell for 2-commutative diagrams.
\xyoption{2cell}
\UseAllTwocells

% We use multicol for the list of chapters between chapters
\usepackage{multicol}

% This is generally recommended for better output
\usepackage{lmodern}
\usepackage[T1]{fontenc}

% For cross-file-references

% Package for hypertext links:
\usepackage{hyperref}

% For any local file, say "hello.tex" you want to link to please

% Theorem environments.
%
\theoremstyle{plain}
\newtheorem{theorem}[subsection]{Theorem}
\newtheorem{proposition}[subsection]{Proposition}
\newtheorem{lemma}[subsection]{Lemma}

\theoremstyle{definition}
\newtheorem{definition}[subsection]{Definition}
\newtheorem{example}[subsection]{Example}
\newtheorem{exercise}[subsection]{Exercise}
\newtheorem{situation}[subsection]{Situation}

\theoremstyle{remark}
\newtheorem{remark}[subsection]{Remark}
\newtheorem{remarks}[subsection]{Remarks}

\numberwithin{equation}{subsection}

% Macros
%
\def\lim{\mathop{\mathrm{lim}}\nolimits}
\def\colim{\mathop{\mathrm{colim}}\nolimits}
\def\Spec{\mathop{\mathrm{Spec}}}
\def\Hom{\mathop{\mathrm{Hom}}\nolimits}
\def\Ext{\mathop{\mathrm{Ext}}\nolimits}
\def\SheafHom{\mathop{\mathcal{H}\!\mathit{om}}\nolimits}
\def\SheafExt{\mathop{\mathcal{E}\!\mathit{xt}}\nolimits}
\def\Sch{\mathit{Sch}}
\def\Mor{\mathop{\mathrm{Mor}}\nolimits}
\def\Ob{\mathop{\mathrm{Ob}}\nolimits}
\def\Sh{\mathop{\mathit{Sh}}\nolimits}
\def\NL{\mathop{N\!L}\nolimits}
\def\CH{\mathop{\mathrm{CH}}\nolimits}
\def\proetale{{pro\text{-}\acute{e}tale}}
\def\etale{{\acute{e}tale}}
\def\QCoh{\mathit{QCoh}}
\def\Ker{\mathop{\mathrm{Ker}}}
\def\Im{\mathop{\mathrm{Im}}}
\def\Coker{\mathop{\mathrm{Coker}}}
\def\Coim{\mathop{\mathrm{Coim}}}

% Boxtimes
%
\DeclareMathSymbol{\boxtimes}{\mathbin}{AMSa}{"02}

%
% Macros for moduli stacks/spaces
%
\def\QCohstack{\mathcal{QC}\!\mathit{oh}}
\def\Cohstack{\mathcal{C}\!\mathit{oh}}
\def\Spacesstack{\mathcal{S}\!\mathit{paces}}
\def\Quotfunctor{\mathrm{Quot}}
\def\Hilbfunctor{\mathrm{Hilb}}
\def\Curvesstack{\mathcal{C}\!\mathit{urves}}
\def\Polarizedstack{\mathcal{P}\!\mathit{olarized}}
\def\Complexesstack{\mathcal{C}\!\mathit{omplexes}}
% \Pic is the operator that assigns to X its picard group, usage \Pic(X)
% \Picardstack_{X/B} denotes the Picard stack of X over B
% \Picardfunctor_{X/B} denotes the Picard functor of X over B
\def\Pic{\mathop{\mathrm{Pic}}\nolimits}
\def\Picardstack{\mathcal{P}\!\mathit{ic}}
\def\Picardfunctor{\mathrm{Pic}}
\def\Deformationcategory{\mathcal{D}\!\mathit{ef}}

\setlength{\emergencystretch}{3em}
\begin{document}
\title[Stacks Project, Tag 095T]{286 --- Proper base change for arbitrary torsion etale coefficients}
\maketitle
\noindent Copyright (C) 2005--2025 Johan de Jong. Stacks Project authors. Source: \href{https://stacks.math.columbia.edu/tag/095T}{Tag 095T}.

\noindent Editorial difficulty note (not author text). Difficulty H3: After a nontrivial reduction to the projective line, one must compute cohomology over a strictly henselian local base using restricted injective resolutions. The proof compares the resulting stalks with closed-fibre cohomology and checks arbitrary field extensions, then promotes prime coefficients to all torsion via the injective criterion..

\noindent Editorial provenance note (not author text). History: excerpt from revision \texttt{a0444\allowbreak 6e57e\allowbreak c1fbc\allowbreak 252a8\allowbreak 71afc\allowbreak ec775\allowbreak 2fb28\allowbreak 07b14}, etale-cohomology.tex, lines 18609--18705. Reference presentation was adapted; original-author context and core support blocks were retained or added. The mathematical wording is unchanged. Licensed under GNU FDL 1.2 or later; no invariant sections or cover texts. See ../licenses/COPYING. Context blocks reproduce author definitions and supporting results; they are not additional counted problems.

\begin{definition}
\label{schemes-definition-reduced-induced-scheme}
Let $X$ be a scheme. Let $Z \subset X$ be a closed subset.
A {\it scheme structure on $Z$} is given by a closed subscheme $Z'$ of
$X$ whose underlying set is equal to $Z$. We often say
``let $(Z, \mathcal{O}_Z)$ be a scheme structure on $Z$'' to
indicate this. The {\it reduced induced scheme structure}
on $Z$ is the one constructed in Lemma \href{https://stacks.math.columbia.edu/tag/01J3}{lemma reduced closed subscheme (Tag 01J3)}.
The {\it reduction $X_{red}$ of $X$} is the reduced induced scheme
structure on $X$ itself.
\end{definition}

\begin{lemma}
\label{lemma-reduce-to-P1}
To prove that cohomology commutes with base change for
every proper morphism of schemes it suffices to prove it
holds for the morphism $\mathbf{P}^1_S \to S$ for every scheme $S$.
\end{lemma}

\begin{proof}
Let $f : X \to Y$ be a proper morphism of schemes.
Let $Y = \bigcup Y_i$ be an affine open covering
and set $X_i = f^{-1}(Y_i)$. If we can prove
cohomology commutes with base change for $X_i \to Y_i$,
then cohomology commutes with base change for $f$.
Namely, the formation of the higher direct images
commutes with Zariski (and even \'etale) localization
on the base, see
Lemma \href{https://stacks.math.columbia.edu/tag/03Q8}{lemma higher direct images (Tag 03Q8)}.
Thus we may assume $Y$ is affine.

\medskip\noindent
Let $Y$ be an affine scheme and let $X \to Y$ be a proper morphism.
By Chow's lemma there exists a commutative diagram
$$
\xymatrix{
X \ar[rd] & X' \ar[d] \ar[l]^\pi \ar[r] & \mathbf{P}^n_Y \ar[dl] \\
& Y &
}
$$
where $X' \to \mathbf{P}^n_Y$ is an immersion, and
$\pi : X' \to X$ is proper and surjective, see
Limits, Lemma \ref{limits-lemma-chow-finite-type}.
Since $X \to Y$ is proper, we find that $X' \to Y$ is proper
(Morphisms, Lemma \href{https://stacks.math.columbia.edu/tag/01W3}{lemma composition proper (Tag 01W3)}).
Hence $X' \to \mathbf{P}^n_Y$ is a closed immersion
(Morphisms, Lemma \href{https://stacks.math.columbia.edu/tag/01W6}{lemma image proper scheme closed (Tag 01W6)}).
It follows that $X' \to X \times_Y \mathbf{P}^n_Y = \mathbf{P}^n_X$
is a closed immersion (as an immersion with closed image).

\medskip\noindent
By Lemma \href{https://stacks.math.columbia.edu/tag/0A4C}{lemma sandwich (Tag 0A4C)}
it suffices to prove cohomology commutes with base change for
$\pi$ and $X' \to Y$. These morphisms both factor as a closed
immersion followed by a projection $\mathbf{P}^n_S \to S$ (for some $S$).
By Lemma \href{https://stacks.math.columbia.edu/tag/0A4E}{lemma finite (Tag 0A4E)} the result holds for closed
immersions (as closed immersions are finite).
By Lemma \href{https://stacks.math.columbia.edu/tag/0A4D}{lemma composition (Tag 0A4D)} it suffices to prove the
result for projections $\mathbf{P}^n_S \to S$.

\medskip\noindent
For every $n \geq 1$ there is a finite surjective morphism
$$
\mathbf{P}^1_S \times_S \ldots \times_S \mathbf{P}^1_S
\longrightarrow
\mathbf{P}^n_S
$$
given on coordinates by
$$
((x_1 : y_1), (x_2 : y_2), \ldots, (x_n : y_n))
\longmapsto
(F_0 : \ldots : F_n)
$$
where $F_0, \ldots, F_n$ in $x_1, \ldots, y_n$
are the polynomials with integer coefficients such that
$$
\prod (x_i t + y_i) = F_0 t^n + F_1 t^{n - 1} + \ldots + F_n
$$
Applying
Lemmas \href{https://stacks.math.columbia.edu/tag/0A4C}{lemma sandwich (Tag 0A4C)}, \href{https://stacks.math.columbia.edu/tag/0A4E}{lemma finite (Tag 0A4E)}, and \href{https://stacks.math.columbia.edu/tag/0A4D}{lemma composition (Tag 0A4D)}
one more time we conclude that the lemma is true.
\end{proof}


\begin{lemma}
\label{lemma-efface-cohomology-on-fibre-by-finite-cover}
Let $A$ be a henselian local ring. Let $X = \mathbf{P}^1_A$.
Let $X_0 \subset X$ be the closed fibre. Let $\ell$ be a prime
number. Let $\mathcal{I}$ be an injective sheaf of
$\mathbf{Z}/\ell\mathbf{Z}$-modules on $X_\etale$. Then
$H^q_\etale(X_0, \mathcal{I}|_{X_0}) = 0$ for $q > 0$.
\end{lemma}

\begin{proof}
Observe that $X$ is a separated scheme which can be covered by $2$
affine opens. Hence for $q > 1$ this follows from Gabber's affine
variant of the proper base change theorem, see
Lemma \ref{lemma-vanishing-restriction-injective}.
Thus we may assume $q = 1$. Let
$\xi \in H^1_\etale(X_0, \mathcal{I}|_{X_0})$.
Goal: show that $\xi$ is $0$.
By Lemmas \href{https://stacks.math.columbia.edu/tag/03SA}{lemma torsion colimit constructible (Tag 03SA)} and
\href{https://stacks.math.columbia.edu/tag/03Q5}{lemma colimit (Tag 03Q5)} we can find a map $\mathcal{F} \to \mathcal{I}$
with $\mathcal{F}$ a constructible sheaf of
$\mathbf{Z}/\ell\mathbf{Z}$-modules
and $\xi$ coming from an element $\zeta$ of
$H^1_\etale(X_0, \mathcal{F}|_{X_0})$. Suppose we have an injective map
$\mathcal{F} \to \mathcal{F}'$ of sheaves of
$\mathbf{Z}/\ell\mathbf{Z}$-modules on $X_\etale$.
Since $\mathcal{I}$ is injective we can extend the given map
$\mathcal{F} \to \mathcal{I}$ to a map $\mathcal{F}' \to \mathcal{I}$.
In this situation we may replace $\mathcal{F}$ by $\mathcal{F}'$
and $\zeta$ by the image of $\zeta$ in $H^1_\etale(X_0, \mathcal{F}'|_{X_0})$.
Also, if $\mathcal{F} = \mathcal{F}_1 \oplus \mathcal{F}_2$ is a direct sum,
then we may replace $\mathcal{F}$ by $\mathcal{F}_i$
and $\zeta$ by the image of $\zeta$ in $H^1_\etale(X_0, \mathcal{F}_i|_{X_0})$.

\medskip\noindent
By Lemma \ref{lemma-constructible-maps-into-constant-general}
and the remarks above we may assume $\mathcal{F}$
is of the form $f_*\underline{M}$ where $M$ is a finite
$\mathbf{Z}/\ell\mathbf{Z}$-module
and $f : Y \to X$ is a finite morphism of finite presentation
(such sheaves are still constructible by
Lemma \href{https://stacks.math.columbia.edu/tag/095R}{lemma finite pushforward constructible (Tag 095R)}
but we won't need this).
Since formation of $f_*$ commutes with any base change
(Lemma \href{https://stacks.math.columbia.edu/tag/0959}{lemma finite pushforward commutes with base change (Tag 0959)})
we see that the restriction of $f_*\underline{M}$ to $X_0$ is
equal to the pushforward of $\underline{M}$ via the induced morphism
$Y_0 \to X_0$ of special fibres. By the Leray spectral sequence
(Proposition \href{https://stacks.math.columbia.edu/tag/03QC}{proposition leray (Tag 03QC)})
and vanishing of higher direct images
(Proposition \href{https://stacks.math.columbia.edu/tag/03QP}{proposition finite higher direct image zero (Tag 03QP)}),
we find
$$
H^1_\etale(X_0, f_*\underline{M}|_{X_0}) = H^1_\etale(Y_0, \underline{M}).
$$
Since $Y \to \Spec(A)$ is proper we can use
Lemma \ref{lemma-proper-over-henselian-and-h1} to see that
the $H^1_\etale(Y_0, \underline{M})$ is equal to
$H^1_\etale(Y, \underline{M})$. Thus we see that our cohomology
class $\zeta$ lifts to a cohomology class
$$
\tilde\zeta \in H^1_\etale(Y, \underline{M}) = H^1_\etale(X, f_*\underline{M})
$$
However, $\tilde \zeta$ maps to zero in
$H^1_\etale(X, \mathcal{I})$ as $\mathcal{I}$ is injective
and by commutativity of
$$
\xymatrix{
H^1_\etale(X, f_*\underline{M}) \ar[r] \ar[d] &
H^1_\etale(X, \mathcal{I}) \ar[d] \\
H^1_\etale(X_0, (f_*\underline{M})|_{X_0}) \ar[r] &
H^1_\etale(X_0, \mathcal{I}|_{X_0})
}
$$
we conclude that the image $\xi$ of $\zeta$ is zero as well.
\end{proof}


\begin{lemma}
\label{lemma-proper-base-change-in-terms-of-injectives}
Let $f : X \to Y$ be a proper morphism of schemes.
The following are equivalent
\begin{enumerate}
\item cohomology commutes with base change for $f$ (see above),
\item for every prime number $\ell$ and every injective
sheaf of $\mathbf{Z}/\ell\mathbf{Z}$-modules $\mathcal{I}$
on $X_\etale$ and every diagram (\ref{equation-base-change-diagram})
where $X' = Y' \times_Y X$ the sheaves
$R^qf'_*(g')^{-1}\mathcal{I}$ are zero for $q > 0$.
\end{enumerate}
\end{lemma}

\begin{proof}
It is clear that (1) implies (2). Conversely, assume (2) and let
$\mathcal{F}$ be a torsion abelian sheaf on $X_\etale$. Let $Y' \to Y$
be a morphism of schemes and let $X' = Y' \times_Y X$
with projections $g' : X' \to X$ and $f' : X' \to Y'$ as in
diagram (\ref{equation-base-change-diagram}).
We want to show the maps of sheaves
$$
g^{-1}R^qf_*\mathcal{F} \longrightarrow R^qf'_*(g')^{-1}\mathcal{F}
$$
are isomorphisms for all $q \geq 0$.

\medskip\noindent
For every $n \geq 1$, let $\mathcal{F}[n]$ be the subsheaf of sections
of $\mathcal{F}$ annihilated by $n$. Then
$\mathcal{F} = \colim \mathcal{F}[n]$.
The functors $g^{-1}$ and $(g')^{-1}$ commute with arbitrary colimits
(as left adjoints). Taking higher direct images along $f$ or $f'$
commutes with filtered colimits by Lemma \href{https://stacks.math.columbia.edu/tag/09Z1}{lemma relative colimit (Tag 09Z1)}.
Hence we see that
$$
g^{-1}R^qf_*\mathcal{F} = \colim g^{-1}R^qf_*\mathcal{F}[n]
\quad\text{and}\quad
R^qf'_*(g')^{-1}\mathcal{F} =
\colim R^qf'_*(g')^{-1}\mathcal{F}[n]
$$
Thus it suffices to prove the result in case $\mathcal{F}$ is
annihilated by a positive integer $n$.

\medskip\noindent
If $n = \ell n'$ for some prime number $\ell$, then we obtain a short
exact sequence
$$
0 \to \mathcal{F}[\ell] \to \mathcal{F} \to
\mathcal{F}/\mathcal{F}[\ell] \to 0
$$
Observe that $\mathcal{F}/\mathcal{F}[\ell]$ is annihilated by $n'$.
Moreover, if the result holds for both $\mathcal{F}[\ell]$ and
$\mathcal{F}/\mathcal{F}[\ell]$, then the result holds by
the long exact sequence of higher direct images (and the $5$ lemma).
In this way we reduce to the case that $\mathcal{F}$ is annihilated
by a prime number $\ell$.

\medskip\noindent
Assume $\mathcal{F}$ is annihilated by a prime number $\ell$.
Choose an injective resolution $\mathcal{F} \to \mathcal{I}^\bullet$
in $D(X_\etale, \mathbf{Z}/\ell\mathbf{Z})$. Applying assumption
(2) and Leray's acyclicity lemma
(Derived Categories, Lemma \href{https://stacks.math.columbia.edu/tag/015E}{lemma leray acyclicity (Tag 015E)})
we see that
$$
f'_*(g')^{-1}\mathcal{I}^\bullet
$$
computes $Rf'_*(g')^{-1}\mathcal{F}$. We conclude by applying
Lemma \ref{lemma-proper-base-change-f-star}.
\end{proof}


\begin{lemma}
\label{lemma-vanishing-restriction-injective}
Let $X$ be a scheme with affine diagonal which can be covered by
$n + 1$ affine opens. Let $Z \subset X$ be a closed subscheme.
Let $\mathcal{A}$ be a torsion sheaf of rings on $X_\etale$
and let $\mathcal{I}$ be an injective sheaf of $\mathcal{A}$-modules
on $X_\etale$.
Then $H^q_\etale(Z, \mathcal{I}|_Z) = 0$ for $q > n$.
\end{lemma}

\begin{proof}
We will prove this by induction on $n$. If $n = 0$, then $X$ is affine.
Say $X = \Spec(A)$ and $Z = \Spec(A/I)$. Let $A^h$ be the filtered colimit
of \'etale $A$-algebras $B$ such that $A/I \to B/IB$ is an isomorphism.
Then $(A^h, IA^h)$ is a henselian pair and $A/I = A^h/IA^h$, see
More on Algebra, Lemma \href{https://stacks.math.columbia.edu/tag/0A02}{lemma henselization (Tag 0A02)}
and its proof. Set $X^h = \Spec(A^h)$.
By Theorem \href{https://stacks.math.columbia.edu/tag/09ZI}{theorem gabber (Tag 09ZI)}
we see that
$$
H^q_\etale(Z, \mathcal{I}|_Z) = H^q_\etale(X^h, \mathcal{I}|_{X^h})
$$
By Theorem \href{https://stacks.math.columbia.edu/tag/09YQ}{theorem colimit (Tag 09YQ)} we have
$$
H^q_\etale(X^h, \mathcal{I}|_{X^h}) =
\colim_{A \to B} H^q_\etale(\Spec(B), \mathcal{I}|_{\Spec(B)})
$$
where the colimit is over the $A$-algebras $B$ as above.
Since the morphisms $\Spec(B) \to \Spec(A)$ are \'etale,
the restriction $\mathcal{I}|_{\Spec(B)}$ is an injective
sheaf of $\mathcal{A}|_{\Spec(B)}$-modules
(Cohomology on Sites, Lemma \href{https://stacks.math.columbia.edu/tag/03F3}{lemma cohomology of open (Tag 03F3)}).
Thus the cohomology groups on the right are zero and we get the
result in this case.

\medskip\noindent
Induction step. We can use Mayer-Vietoris to do the induction step.
Namely, suppose that $X = U \cup V$ where $U$ is a union of $n$ affine
opens and $V$ is affine. Then, using that the diagonal of $X$ is affine,
we see that $U \cap V$ is the union of $n$ affine opens. Mayer-Vietoris
gives an exact sequence
$$
H^{q - 1}_\etale(U \cap V \cap Z, \mathcal{I}|_Z) \to
H^q_\etale(Z, \mathcal{I}|_Z) \to
H^q_\etale(U \cap Z, \mathcal{I}|_Z) \oplus
H^q_\etale(V \cap Z, \mathcal{I}|_Z)
$$
and by our induction hypothesis we obtain vanishing for $q > n$ as desired.
\end{proof}


\begin{lemma}
\label{lemma-constructible-maps-into-constant-general}
\begin{reference}
\href{https://stacks.math.columbia.edu/bibliography}{Bibliography: SGA4 (Exposee IX, Proposition 2.14)}
\end{reference}
Let $X$ be a quasi-compact and quasi-separated scheme.
\begin{enumerate}
\item Let $\mathcal{F}$ be a constructible sheaf of sets on $X_\etale$.
There exist an injective map of sheaves
$$
\mathcal{F} \longrightarrow
\prod\nolimits_{i = 1, \ldots, n} f_{i, *}\underline{E_i}
$$
where $f_i : Y_i \to X$ is a finite and finitely presented morphism and
$E_i$ is a finite set.
\item Let $\mathcal{F}$ be a constructible abelian sheaf on $X_\etale$.
There exist an injective map of abelian sheaves
$$
\mathcal{F} \longrightarrow
\bigoplus\nolimits_{i = 1, \ldots, n} f_{i, *}\underline{M_i}
$$
where $f_i : Y_i \to X$ is a finite and finitely presented morphism and
$M_i$ is a finite abelian group.
\item Let $\Lambda$ be a Noetherian ring.
Let $\mathcal{F}$ be a constructible sheaf of $\Lambda$-modules on $X_\etale$.
There exist an injective map of sheaves of modules
$$
\mathcal{F} \longrightarrow
\bigoplus\nolimits_{i = 1, \ldots, n} f_{i, *}\underline{M_i}
$$
where $f_i : Y_i \to X$ is a finite and finitely presented morphism and
$M_i$ is a finite $\Lambda$-module.
\end{enumerate}
\end{lemma}

\begin{proof}
We will reduce this lemma to the Noetherian case by absolute Noetherian
approximation. Namely, by
Limits, Proposition \href{https://stacks.math.columbia.edu/tag/01ZA}{proposition approximate (Tag 01ZA)}
we can write $X = \lim_{t \in T} X_t$ with each $X_t$ of finite type over
$\Spec(\mathbf{Z})$ and with affine transition morphisms. By
Lemma \href{https://stacks.math.columbia.edu/tag/09YU}{lemma category is colimit (Tag 09YU)}
the category of constructible sheaves (of sets, abelian groups, or
$\Lambda$-modules) on $X_\etale$ is the colimit of the corresponding
categories for $X_t$. Thus our constructible sheaf $\mathcal{F}$
is the pullback of a similar constructible sheaf $\mathcal{F}_t$
over $X_t$ for some $t$. Then we apply the Noetherian case
(Lemma \ref{lemma-constructible-maps-into-constant})
to find an injection
$$
\mathcal{F}_t \longrightarrow
\prod\nolimits_{i = 1, \ldots, n} f_{i, *}\underline{E_i}
\quad\text{or}\quad
\mathcal{F}_t \longrightarrow
\bigoplus\nolimits_{i = 1, \ldots, n} f_{i, *}\underline{M_i}
$$
over $X_t$ for some finite morphisms $f_i : Y_i \to X_t$.
Since $X_t$ is Noetherian the morphisms $f_i$ are of finite presentation.
Since pullback is exact and since formation of $f_{i, *}$ commutes
with base change
(Lemma \href{https://stacks.math.columbia.edu/tag/0959}{lemma finite pushforward commutes with base change (Tag 0959)}), we conclude.
\end{proof}


\begin{lemma}
\label{lemma-constructible-maps-into-constant}
Let $X$ be a Noetherian scheme.
\begin{enumerate}
\item Let $\mathcal{F}$ be a constructible sheaf of sets on $X_\etale$.
There exist an injective map of sheaves
$$
\mathcal{F} \longrightarrow
\prod\nolimits_{i = 1, \ldots, n} f_{i, *}\underline{E_i}
$$
where $f_i : Y_i \to X$ is a finite morphism and $E_i$ is a finite set.
\item Let $\mathcal{F}$ be a constructible abelian sheaf on $X_\etale$.
There exist an injective map of abelian sheaves
$$
\mathcal{F} \longrightarrow
\bigoplus\nolimits_{i = 1, \ldots, n} f_{i, *}\underline{M_i}
$$
where $f_i : Y_i \to X$ is a finite morphism and
$M_i$ is a finite abelian group.
\item Let $\Lambda$ be a Noetherian ring.
Let $\mathcal{F}$ be a constructible sheaf of $\Lambda$-modules on $X_\etale$.
There exist an injective map of sheaves of modules
$$
\mathcal{F} \longrightarrow
\bigoplus\nolimits_{i = 1, \ldots, n} f_{i, *}\underline{M_i}
$$
where $f_i : Y_i \to X$ is a finite morphism and
$M_i$ is a finite $\Lambda$-module.
\end{enumerate}
Moreover, we may assume each $Y_i$ is irreducible, reduced, maps onto
an irreducible and reduced closed subscheme $Z_i \subset X$ such that
$Y_i \to Z_i$ is finite \'etale over a nonempty open of $Z_i$.
\end{lemma}

\begin{proof}
Proof of (1). Because we have the ascending chain condition for
subsheaves of $\mathcal{F}$
(Lemma \href{https://stacks.math.columbia.edu/tag/09YV}{lemma constructible over noetherian noetherian (Tag 09YV)}), it
suffices to show that for every point $x \in X$ we
can find a map $\varphi : \mathcal{F} \to f_*\underline{E}$ where
$f : Y \to X$ is finite and $E$ is a finite set such that
$\varphi_{\overline{x}} : \mathcal{F}_{\overline{x}} \to
(f_*S)_{\overline{x}}$ is injective.
(This argument can be avoided by picking a partition of $X$ as in
Lemma \href{https://stacks.math.columbia.edu/tag/095E}{lemma constructible quasi compact quasi separated (Tag 095E)}
and constructing a $Y_i \to X$ for each irreducible component
of each part.)
Let $Z \subset X$ be the induced reduced scheme structure
(Schemes, Definition \ref{schemes-definition-reduced-induced-scheme})
on $\overline{\{x\}}$.
Since $\mathcal{F}$ is constructible, there is a finite separable
extension $K/\kappa(x)$ such that
$\mathcal{F}|_{\Spec(K)}$ is the constant sheaf with value $E$
for some finite set $E$. Let $Y \to Z$ be the normalization
of $Z$ in $\Spec(K)$.
By Morphisms, Lemma \href{https://stacks.math.columbia.edu/tag/035L}{lemma normal normalization (Tag 035L)}
we see that $Y$ is a normal integral scheme.
As $K/\kappa(x)$ is a finite extension, it is clear that $K$ is the function
field of $Y$. Denote $g : \Spec(K) \to Y$ the inclusion.
The map $\mathcal{F}|_{\Spec(K)} \to \underline{E}$ is adjoint
to a map $\mathcal{F}|_Y \to g_*\underline{E} = \underline{E}$
(Lemma \href{https://stacks.math.columbia.edu/tag/09Z5}{lemma push constant sheaf from open (Tag 09Z5)}).
This in turn is adjoint to a map
$\varphi : \mathcal{F} \to f_*\underline{E}$.
Observe that the stalk of $\varphi$ at a geometric point
$\overline{x}$ is injective: we may take a lift $\overline{y} \in Y$
of $\overline{x}$ and the commutative diagram
$$
\xymatrix{
\mathcal{F}_{\overline{x}} \ar@{=}[r] \ar[d] &
(\mathcal{F}|_Y)_{\overline{y}} \ar@{=}[d] \\
(f_*\underline{E})_{\overline{x}} \ar[r] &
\underline{E}_{\overline{y}}
}
$$
proves the injectivity. We are not yet done, however, as the
morphism $f : Y \to Z$ is integral but in general not
finite\footnote{If $X$ is a Nagata scheme, for example of finite
type over a field, then $Y \to Z$ is finite.}.

\medskip\noindent
To fix the problem stated in the last sentence of the previous paragraph,
we write $Y = \lim_{i \in I} Y_i$ with $Y_i$ irreducible, integral, and
finite over $Z$. Namely, apply Properties, Lemma
\href{https://stacks.math.columbia.edu/tag/0817}{lemma integral algebra directed colimit finite (Tag 0817)}
to $f_*\mathcal{O}_Y$ viewed as a sheaf of $\mathcal{O}_Z$-algebras
and apply the functor $\underline{\Spec}_Z$.
Then $f_*\underline{E} = \colim f_{i, *}\underline{E}$
by Lemma \href{https://stacks.math.columbia.edu/tag/09Z1}{lemma relative colimit (Tag 09Z1)}.
By Lemma \href{https://stacks.math.columbia.edu/tag/09Z4}{lemma constructible is compact (Tag 09Z4)} the map
$\mathcal{F} \to f_*\underline{E}$
factors through $f_{i, *}\underline{E}$ for some $i$.
Since $Y_i \to Z$ is a finite morphism of integral schemes
and since the function field extension
induced by this morphism is finite separable, we see that the
morphism is finite \'etale over a nonempty open of $Z$ (use
Algebra, Lemma \href{https://stacks.math.columbia.edu/tag/07ND}{lemma smooth at generic point (Tag 07ND)}; details omitted).
This finishes the proof of (1).

\medskip\noindent
The proofs of (2) and (3) are identical to the proof of (1).
\end{proof}


\begin{lemma}
\label{lemma-proper-over-henselian-and-h1}
Let $A$ be a henselian local ring. Let $X$ be a proper scheme over $A$
with closed fibre $X_0$. Let $M$ be a finite abelian group.
Then $H^1_\etale(X, \underline{M}) = H^1_\etale(X_0, \underline{M})$.
\end{lemma}

\begin{proof}
By Cohomology on Sites, Lemma \href{https://stacks.math.columbia.edu/tag/03AJ}{lemma torsors h1 (Tag 03AJ)}
an element of $H^1_\etale(X, \underline{M})$ corresponds to a
$\underline{M}$-torsor $\mathcal{F}$ on $X_\etale$.
Such a torsor is clearly a finite locally constant sheaf.
Hence $\mathcal{F}$ is representable by a scheme $V$ finite
\'etale over $X$, Lemma \href{https://stacks.math.columbia.edu/tag/03RV}{lemma characterize finite locally constant (Tag 03RV)}.
Conversely, a scheme $V$ finite \'etale over $X$ with an $M$-action
which turns it into an $M$-torsor over $X$ gives rise to a cohomology
class. The same translation between cohomology classes over $X_0$ and
torsors finite \'etale over $X_0$ holds. Thus the lemma
is a consequence of the equivalence of categories of
Fundamental Groups, Lemma
\href{https://stacks.math.columbia.edu/tag/0A48}{lemma finite etale on proper over henselian (Tag 0A48)}.
\end{proof}


\begin{lemma}
\label{lemma-proper-base-change-f-star}
Let $f : X \to Y$ be a proper morphism of schemes. Let $g : Y' \to Y$
be a morphism of schemes. Set $X' = Y' \times_Y X$ with projections
$f' : X' \to Y'$ and $g' : X' \to X$. Let $\mathcal{F}$ be any sheaf on
$X_\etale$. Then $g^{-1}f_*\mathcal{F} = f'_*(g')^{-1}\mathcal{F}$.
\end{lemma}

\begin{proof}
There is a canonical map $g^{-1}f_*\mathcal{F} \to f'_*(g')^{-1}\mathcal{F}$.
Namely, it is adjoint to the map
$$
f_*\mathcal{F} \longrightarrow
g_*f'_*(g')^{-1}\mathcal{F} = f_*g'_*(g')^{-1}\mathcal{F}
$$
which is $f_*$ applied to the canonical map
$\mathcal{F} \to g'_*(g')^{-1}\mathcal{F}$. To check this map is an
isomorphism we can compute what happens on stalks.
Let $y' : \Spec(k) \to Y'$ be a geometric point with image $y$ in $Y$.
By Lemma \href{https://stacks.math.columbia.edu/tag/0A3T}{lemma proper pushforward stalk (Tag 0A3T)} the stalks are
$\Gamma(X'_{y'}, \mathcal{F}_{y'})$ and $\Gamma(X_y, \mathcal{F}_y)$
respectively. Here the sheaves $\mathcal{F}_y$ and $\mathcal{F}_{y'}$
are the pullbacks of $\mathcal{F}$ by the projections $X_y \to X$
and $X'_{y'} \to X$. Since $X_y = X'_{y'}$ as schemes
we see that the stalks agree. We omit the
verification that this isomorphism is compatible with our map.
\end{proof}


\begin{lemma}
\label{limits-lemma-chow-finite-type}
Let $S$ be a quasi-compact and quasi-separated scheme.
Let $f : X \to S$ be a separated morphism of finite type.
Then there exists an $n \geq 0$ and a diagram
$$
\xymatrix{
X \ar[rd] & X' \ar[d] \ar[l]^\pi \ar[r] & \mathbf{P}^n_S \ar[dl] \\
& S &
}
$$
where $X' \to \mathbf{P}^n_S$ is an immersion, and
$\pi : X' \to X$ is proper and surjective.
\end{lemma}

\begin{proof}
By Proposition \href{https://stacks.math.columbia.edu/tag/01ZJ}{proposition separated closed in finite presentation (Tag 01ZJ)}
we can find a closed immersion $X \to Y$ where $Y$ is separated
and of finite presentation over $S$. Clearly, if we prove the assertion
for $Y$, then the result follows for $X$. Hence we may assume that
$X$ is of finite presentation over $S$.

\medskip\noindent
Write $S = \lim_i S_i$ as a directed limit of Noetherian schemes, see
Proposition \href{https://stacks.math.columbia.edu/tag/01ZA}{proposition approximate (Tag 01ZA)}. By
Lemma \href{https://stacks.math.columbia.edu/tag/01ZM}{lemma descend finite presentation (Tag 01ZM)} we can
find an index $i \in I$ and a scheme $X_i \to S_i$ of finite presentation
so that $X = S \times_{S_i} X_i$.
By Lemma \href{https://stacks.math.columbia.edu/tag/01ZQ}{lemma descend separated finite presentation (Tag 01ZQ)}
we may assume that $X_i \to S_i$ is separated.
Clearly, if we prove the assertion for
$X_i$ over $S_i$, then the assertion holds for $X$. The case
$X_i \to S_i$ is treated by
Cohomology of Schemes, Lemma \href{https://stacks.math.columbia.edu/tag/0200}{lemma chow Noetherian (Tag 0200)}.
\end{proof}


\noindent
At this point we start discussing the proper base change theorem.
To do so we introduce some notation. consider a commutative diagram
\begin{equation}
\label{equation-base-change-diagram}
\vcenter{
\xymatrix{
X' \ar[r]_{g'} \ar[d]_{f'} & X \ar[d]^f \\
Y' \ar[r]^g & Y
}
}
\end{equation}
of morphisms of schemes. Then we obtain a commutative diagram of sites
$$
\xymatrix{
X'_\etale \ar[r]_{g'_{small}} \ar[d]_{f'_{small}} &
X_\etale \ar[d]^{f_{small}} \\
Y'_\etale \ar[r]^{g_{small}} &
Y_\etale
}
$$
For any object $E$ of $D(X_\etale)$ we obtain a canonical base change map
\begin{equation}
\label{equation-base-change}
g_{small}^{-1}Rf_{small, *}E \longrightarrow Rf'_{small, *}(g'_{small})^{-1}E
\end{equation}
in $D(Y'_\etale)$. See Cohomology on Sites, Remark
\href{https://stacks.math.columbia.edu/tag/07A7}{remark base change (Tag 07A7)} where we use the constant
sheaf $\mathbf{Z}$ as our sheaf of rings.
We will usually omit the subscripts ${}_{small}$ in this formula.
For example, if $E = \mathcal{F}[0]$ where $\mathcal{F}$ is an abelian
sheaf on $X_\etale$, the base change map is a map
\begin{equation}
\label{equation-base-change-sheaf}
g^{-1}Rf_*\mathcal{F} \longrightarrow Rf'_*(g')^{-1}\mathcal{F}
\end{equation}
in $D(Y'_\etale)$.

\medskip\noindent
The map (\ref{equation-base-change}) has no chance of being an isomorphism
in the generality given above. The goal is to show
it is an isomorphism if the diagram (\ref{equation-base-change-diagram})
is cartesian, $f : X \to Y$ proper, the cohomology sheaves of $E$
are torsion, and $E$ is bounded below.
To study this question we introduce the following terminology.
Let us say that {\it cohomology commutes with base change
for $f : X \to Y$} if (\ref{equation-base-change-sheaf})
is an isomorphism for every diagram (\ref{equation-base-change-diagram})
where $X' = Y' \times_Y X$ and every torsion abelian sheaf $\mathcal{F}$.


\begin{theorem}[Proper Base Change]
\label{theorem-proper-base-change}
Let $f : X \to Y$ be a proper morphism of schemes. Let $g : Y' \to Y$ be
a morphism of schemes. Set $X' = Y' \times_Y X$
and consider the cartesian diagram
$$
\xymatrix{
X' \ar[r]_{g'} \ar[d]_{f'} & X \ar[d]^f \\
Y' \ar[r]^g & Y
}
$$
Let $\mathcal{F}$ be an abelian torsion sheaf on $X_\etale$.
Then the base change map
$$
g^{-1}Rf_*\mathcal{F} \longrightarrow Rf'_*(g')^{-1}\mathcal{F}
$$
is an isomorphism.
\end{theorem}

\begin{proof}
In the terminology introduced above, this means that cohomology commutes
with base change for every proper morphism of schemes. By
Lemma \ref{lemma-reduce-to-P1}
it suffices to prove that cohomology commutes with base change
for the morphism $\mathbf{P}^1_S \to S$ for every scheme $S$.

\medskip\noindent
Let $S$ be the spectrum of a strictly henselian local ring with closed
point $s$. Set $X = \mathbf{P}^1_S$ and $X_0 = X_s = \mathbf{P}^1_s$.
Let $\mathcal{F}$ be a sheaf of $\mathbf{Z}/\ell\mathbf{Z}$-modules
on $X_\etale$. The key to our proof is that
$$
H^q_\etale(X, \mathcal{F}) = H^q_\etale(X_0, \mathcal{F}|_{X_0}).
$$
Namely, choose a resolution $\mathcal{F} \to \mathcal{I}^\bullet$
by injective sheaves of $\mathbf{Z}/\ell\mathbf{Z}$-modules.
Then $\mathcal{I}^\bullet|_{X_0}$ is a resolution of $\mathcal{F}|_{X_0}$
by right $H^0_\etale(X_0, -)$-acyclic objects, see
Lemma \ref{lemma-efface-cohomology-on-fibre-by-finite-cover}.
Leray's acyclicity lemma tells us the right hand side is computed by
the complex $H^0_\etale(X_0, \mathcal{I}^\bullet|_{X_0})$
which is equal to $H^0_\etale(X, \mathcal{I}^\bullet)$ by
Lemma \href{https://stacks.math.columbia.edu/tag/0A3S}{lemma h0 proper over henselian local (Tag 0A3S)}. This complex
computes the left hand side.

\medskip\noindent
Assume $S$ is general and $\mathcal{F}$ is a sheaf of
$\mathbf{Z}/\ell\mathbf{Z}$-modules on $X_\etale$.
Let $\overline{s} : \Spec(k) \to S$ be a geometric point
of $S$ lying over $s \in S$. We have
$$
(R^qf_*\mathcal{F})_{\overline{s}} =
H^q_\etale(\mathbf{P}^1_{\mathcal{O}_{S, \overline{s}}^{sh}},
\mathcal{F}|_{\mathbf{P}^1_{\mathcal{O}_{S, \overline{s}}^{sh}}}) =
H^q_\etale(\mathbf{P}^1_{\kappa(s)^{sep}},
\mathcal{F}|_{\mathbf{P}^1_{\kappa(s)^{sep}}})
$$
where $\kappa(s)^{sep}$ is the residue field of
$\mathcal{O}_{S, \overline{s}}^{sh}$, i.e., the separable algebraic
closure of $\kappa(s)$ in $k$.
The first equality by Theorem \href{https://stacks.math.columbia.edu/tag/03Q9}{theorem higher direct images (Tag 03Q9)}
and the second equality by the displayed formula in the
previous paragraph.

\medskip\noindent
Finally, consider any morphism of schemes $g : T \to S$ where
$S$ and $\mathcal{F}$ are as above.
Set $f' : \mathbf{P}^1_T \to T$ the projection and let
$g' : \mathbf{P}^1_T \to \mathbf{P}^1_S$ the morphism induced
by $g$. Consider the base change map
$$
g^{-1}R^qf_*\mathcal{F}
\longrightarrow
R^qf'_*(g')^{-1}\mathcal{F}
$$
Let $\overline{t}$ be a geometric point of $T$ with image
$\overline{s} = g(\overline{t})$. By our discussion
above the map on stalks at $\overline{t}$ is the map
$$
H^q_\etale(\mathbf{P}^1_{\kappa(s)^{sep}},
\mathcal{F}|_{\mathbf{P}^1_{\kappa(s)^{sep}}})
\longrightarrow
H^q_\etale(\mathbf{P}^1_{\kappa(t)^{sep}},
\mathcal{F}|_{\mathbf{P}^1_{\kappa(t)^{sep}}})
$$
Since $\kappa(s)^{sep} \subset \kappa(t)^{sep}$ this map is an
isomorphism by Lemma \href{https://stacks.math.columbia.edu/tag/0A5E}{lemma base change dim 1 separably closed (Tag 0A5E)}.

\medskip\noindent
This proves cohomology commutes with base change for
$\mathbf{P}^1_S \to S$ and sheaves of $\mathbf{Z}/\ell\mathbf{Z}$-modules.
In particular, for an injective sheaf of $\mathbf{Z}/\ell\mathbf{Z}$-modules
the higher direct images of any base change are zero.
In other words, condition (2) of
Lemma \ref{lemma-proper-base-change-in-terms-of-injectives}
holds and the proof is complete.
\end{proof}
\end{document}
